\section{Scope, contributions, and non-claims}
\label{sec:scope}

This paper is organized as a consistency architecture for a deterministic substrate hypothesis.
Its purpose is to identify which connections between the minimal substrate and the observed structure of physics
are already mathematically fixed, which require further derivation, and which must be decided by computation.
A cubic lattice is taken as the candidate microphysical substrate; the continuum formulation serves as
the long-wavelength effective description and the cleanest analytical baseline.

\subsection{Positioning relative to established physics}
The Standard Model and general relativity are treated here as empirically validated effective frameworks.
Nothing in this paper is intended to weaken confidence in their experimental domain of applicability.
The present work instead targets a different explanatory layer: it asks whether familiar structures such as
effective Lorentz symmetry, gauge connections, and quantized particle spectra could arise as emergent regularities
from a smaller set of ground ingredients.

\subsection{Consistency bridges}
\label{subsec:consistency-bridges}
The theoretical framework is organized around six coupled consistency bridges --- each connecting the minimal substrate to a structure that any viable physical ontology must possess:
\begin{enumerate}
\item \textbf{A single substrate with two descriptions.}
A 4D cubic lattice embedded with codimension zero in a 4D Euclidean ambient is taken as the microphysical substrate; the timelike direction is selected by the Lorentzian sign structure of the brane action (\Cref{sec:continuum-model}).
For wavelengths large compared to the lattice spacing, its dynamics on a spacelike slice admits a continuum approximation on an intrinsic $3$-domain $\Omega\subset\R^3$. An isotropic hyperelastic action remains the analytic
baseline, but the current cubic stencil has a measurable leading-order anisotropy unless its shell weights are retuned.

\item \textbf{Geometric nonlinearity as the primary coupling mechanism.}
The central amplitude--lateral coupling does not come from an auxiliary field.
It comes from the 4D Euclidean link length on the lattice and, equivalently, from the induced metric in the continuum.
This geometric mechanism is the intended source of transverse contraction, self-guidance, and localization.

\item \textbf{A prestretched substrate with an emergent gravity-like channel.}
A homogeneous mismatch parameter $\alpha$ encodes the prestretch between held spacing and stress-free length.
Because transverse gradients modify intrinsic distances, localized transverse excitation produces an inward lateral pull
(contraction) and a slowly varying strain field.
In a weak, slowly varying regime this is interpreted as the qualitative origin of a long-range gravity-like channel
(\Cref{subsec:gravity-channel}).

\item \textbf{Per-wavepacket band isolation and complex-envelope promotion.}
The operational ingredient is not a globally coherent carrier sector but per-wavepacket band isolation:
each excitation is treated as narrowband around its own local carrier frequency $\omega_0$, with its own
retained polarization subspace defined as an eigenbundle of the linearized carrier operator about the
slowly varying background at the wavepacket's location.
This local band isolation enables the controlled multiscale envelope description and promotes the real
lateral triplet to $\Psi\in\C^3$ (\Cref{sec:geophase,sec:effective-field}).
No universe-wide coherent carrier is postulated.

\item \textbf{Berry/Wilczek--Zee transport as emergent gauge structure.}
Adiabatic transport of a retained polarization subspace induces a Berry connection (rank-1) or a Wilczek--Zee
connection (non-Abelian).
These are interpreted as the natural gauge variables of the slow-sector description, rather than as independent
microscopic fields (\Cref{sec:geophase,sec:effective-field}).

\item \textbf{Localized particle modes and the baryon priority.}
Particle-like states are modeled as localized standing-wave/solitonic modes whose spatial structure is organized by
symmetry and whose amplitude scale is fixed by nonlinear self-guidance and energy normalization (\Cref{sec:localized}).
Because the cubic substrate naturally singles out an axis-aligned triplet sector, the first decisive non-electron target
is argued to be a baryon-like confined triplet mode, with proton/neutron structure treated as the primary numerical
priority.
\end{enumerate}

\subsection{Non-claims}
To keep the scope honest, several items are explicitly \emph{not} claimed in this manuscript:
\begin{itemize}
\item a derivation of the full Standard Model gauge group, generations, or precise coupling constants,
\item a completed quantitative reduction of the contraction channel to a Newtonian limit with a derived value of $G$,
\item an identification of any specific topological object (vortex ring, texture, etc.) with a specific physical
particle: \Cref{subsec:triplet-u3} records the topological structure of each sector as a mathematical fact, but
which objects the substrate dynamics stabilizes is open,
\item a derivation of the QCD-vs-EM scale hierarchy from the prestress parameter $\alpha$: $\alpha$ controls a
coupling ratio in the linear spectrum but whether a true dynamical scale hierarchy emerges from the nonlinear sector
is unsettled and not claimed here.
\end{itemize}

\subsection{First-layer objections addressed}
\label{subsec:attack-responses}

A substrate hypothesis of this breadth encounters five standard foundational objections.
This paper addresses each as a consistency requirement.

\begin{enumerate}
\item \emph{Lattice breaks Lorentz symmetry.}
A cubic UV substrate is not rotationally symmetric.
The response is the dual-observer argument (\Cref{para:dual-observer}): if both the medium and the inside observers built from it are governed by the same stiffness tensor, the effective kinematics experienced internally may be isotropic even when the lab metric is not.
Residual birefringence is then a falsifiable signature, not a tunable parameter.

\item \emph{Bell's theorem rules out local determinism.}
The substrate is deterministic but not locally separable: entanglement is the continuity of a single 4D worldtube object that branches at a V-vertex in the past.
Causality emerges from soliton chirality, not from the substrate.
This is a retrocausal completion (\Cref{sec:bell-interpretation}), consistent with Bell's inequalities.

\item \emph{Where are the particles?}
No completed particle solution is claimed.
The program is posed as a mathematical eigenproblem (\Cref{sec:localized}): given the substrate action and the geometric self-guidance mechanism, do stable localized standing-wave solutions exist?
The baryon-like triplet simulation is the first decisive test.

\item \emph{Gauge structure as coincidence.}
Berry and Wilczek--Zee connections are not coincidentally gauge-like; they are the canonical geometric structure that must appear whenever a multi-component eigenbundle is adiabatically transported along a path in parameter space (\Cref{sec:geophase,sec:effective-field}).
The present claim is that band-isolated wavepacket eigenbundles are the physically appropriate transport objects.

\item \emph{Probability without fundamentals.}
Apparent probabilistic outcomes are attributed to Fourier-localization constraints and nonlinear exchange events in the deterministic dynamics, not to fundamental stochasticity.
Apparent discreteness reflects wave-packet topology, not fundamental point particles.
\end{enumerate}

\subsection{Immediate numerical program}
The paper now commits to a concrete simulation agenda rather than to a broad list of generic falsification slogans.
Three numerical tasks are decisive:
\begin{enumerate}
\item \textbf{Dispersion and isotropy control.}
Measure direction-dependent wave speeds, birefringence, and branch splitting as functions of wavelength relative to the
lattice spacing and neighbor stencil.
The aim is to determine whether a universal long-wavelength observer sector exists.

\item \textbf{Contraction-channel extraction.}
Create localized transverse excitation, measure the induced slow in-brane contraction field, and test whether it focuses
or accelerates long-wavelength probe packets in the expected weak-field manner.

\item \textbf{Localized triplet-mode search.}
Construct and evolve confined three-component seeds tied to the axis-aligned cubic channels, and test whether the
substrate supports stable baryon-like standing-wave modes with confined internal triplet structure, controlled far-field
holonomy, and reproducible size/energy diagnostics.
\end{enumerate}
The contribution of the present paper is therefore a sharpened substrate model plus a concrete route for deciding whether
its main emergence claims survive direct computation.
